\documentclass[11pt]{article}
\usepackage[margin=1in]{geometry}
\usepackage{amsmath,amssymb,amsthm,mathtools}
\usepackage{hyperref}
\usepackage{enumitem}
\usepackage{booktabs}
\usepackage{array}

\newtheorem{theorem}{Theorem}[section]
\newtheorem{corollary}[theorem]{Corollary}
\newtheorem{lemma}[theorem]{Lemma}
\newtheorem*{definition}{Definition}
\newtheorem*{remark}{Remark}

\newcommand{\Q}{\mathbb{Q}}
\newcommand{\Z}{\mathbb{Z}}
\newcommand{\N}{\mathbb{N}}
\newcommand{\repo}[1]{\text{\texttt{#1}}}
\newcommand{\cw}{c_W}

\hypersetup{colorlinks=true, linkcolor=blue, citecolor=blue, urlcolor=blue}

\title{Beyond Sturmian Words in the 3x+1 Periodicity Conjecture:\\
A Periodic-Approximation and Arithmetic-Height Bridge}
\author{Elias De Jesús\\
\textit{Independent Researcher}\\
ORCID: \href{https://orcid.org/0009-0007-0190-9143}{0009-0007-0190-9143}}
\date{September 2026 \\[2pt] \textnormal{\small Revision 2, 30 September 2026 --- corrected; see \S\ref{sec:rev2}}}

\begin{document}
\sloppy
\maketitle

\begin{abstract}
Let $T:\Z_2\to\Z_2$ be the $3x+1$ map and $\Phi$ the Bernstein--Lagarias conjugacy map sending a parity vector $v\in\{0,1\}^\N$ to the 2-adic integer it codes. Lagarias's Periodicity Conjecture asserts that $\Phi(v)\in\Q$ only if $v$ is eventually periodic; it is open in general. A companion unpublished paper by the present author proves it for every Sturmian parity word, by two independent mechanisms (a transported counting bound, and an effective Liouville-style argument using an initial square). This note consolidates a subsequent research program that pushes the periodic-approximation/arithmetic-height mechanism of the second route past the Sturmian case, into symbolic classes the counting route cannot reach. The main proved result: for \emph{every} irrational slope with bounded partial quotients, every complementary symmetric Rote sequence (factor complexity exactly $2n$) built on \emph{any} Sturmian word of that slope --- any intercept, either mechanical convention, either seed, any shift --- has irrational $\Phi$-value. We give the full proof architecture, including a structural fact (an odd-weight transfer route's doubled root always has density exactly $1/2$) that collapses an earlier, over-conservative repetition-exponent requirement down to the true one. We document, rather than hide, four working hypotheses this program formed and then disproved by direct adversarial testing. A late-stage exploratory probe into a conjectural symbolic-symmetry/arithmetic-admissibility tradeoff is reported honestly as inconclusive, not as a result.

\smallskip
\noindent\textbf{Revision 2.} Revision 1 presented the general-intercept case as blocked by one unproved margin lemma (its Conjecture 12.1). That lemma is \emph{not} open: it is Berth\'e--Holton--Zamboni (2006), Propositions 5.1 and 5.2 --- Section 5 of the same paper this program already cites as its primary source. Revision 2 therefore states the headline theorem at \emph{every} intercept of every bounded-type slope, under either mechanical convention, and corrects four further defects of Revision 1: a mis-indexed convergent-denominator convention, a false ``exceptional-orbit'' reduction, a missing length factor in the surplus decomposition, and an overstated attainment claim for one of the two Ostrowski repetition expressions. The symbolic margin is credited to Berth\'e--Holton--Zamboni throughout; what is original here is the arithmetic application --- the periodic-approximation/height bridge and the attained-witness bookkeeping that converts a $\limsup$ bound on the initial critical exponent into an unbounded surplus. \S\ref{sec:rev2} lists every change and says which Revision-1 conclusions survive; the full audit trail is \cite{Phase8Target,Phase8Verdict}. \textbf{This note does not prove, and does not claim progress toward proving, the Periodicity Conjecture in general}; a structural fact recorded in the underlying research record shows no purely repetition-based method, however extended, can reach every aperiodic word.
\end{abstract}

\tableofcontents

\section{Revision 2: Corrections to Revision 1}
\label{sec:rev2}

Revision 1 of this note (Zenodo version DOI \texttt{10.5281/zenodo.23045141}) is superseded in five
places. All five were found by auditing this program's own record against
Berth\'e--Holton--Zamboni \cite{BHZ2006} read first-hand, including its \S5, which the underlying
Phase 7 work did not reach.

\begin{description}[leftmargin=2.5em,style=nextline]

\item[C1 --- The ``one remaining lemma'' is a published theorem.] Revision 1 \S12 stated as
\emph{Conjecture 12.1} the claim that for every admissible Ostrowski digit sequence of a
bounded-type slope, $\limsup_k\max(x'(k),y(k))>2$ with a margin depending only on $A$, and said it
was ``not claimed proved here or anywhere in the underlying record'' \cite{Phase7Sufficient,Phase7Verdict}. This is \textbf{BHZ
Propositions 5.1 and 5.2} \cite{BHZ2006}: if $(a_k,a_{k+1})=(s,t)$ with $s>1$ for infinitely many
$k$ then \emph{every} $\omega\in X_\alpha$ has $\text{ice}(\omega)\ge2+\frac1{2(s+1)(t+1)+1}$, and if
$(a_k,a_{k+1},a_{k+2})=(1,1,t)$ for infinitely many $k$ then every $\omega\in X_\alpha$ has
$\text{ice}(\omega)\ge2+\frac1{8t+1}$. A bounded-type slope satisfies one hypothesis or the other by
pigeonhole. \emph{Consequence:} Theorem \ref{thm:boundedtype} now holds at every intercept
(\S\ref{sec:generalintercept}), and Revision 1's Open 1 is not open. The weaker statement
$\text{ice}(\omega)>2$ at each individual intercept is likewise already a corollary of BHZ
Theorem 1.1.

\item[C2 --- The ``exceptional-orbit reduction'' is false.] Revision 1's Open 2, and the Phase 7
verdict it quoted, rested on the claim that $\text{ice}(\omega_\rho)=\text{ind}^*(\alpha)$ for every
intercept outside the countable ``return'' set $\{j\alpha\bmod1\}$ \cite{BHZGeneralMap}, so that the genuinely open cases
were countable and intercept $0$ was ``if anything a harder case than generic.'' The cited
ingredients do not give this: BHZ Proposition 2.1(4) gives shift-invariance of $\text{ice}$ only off
a \emph{finite set of orbits}, and Lemma 2.2 gives $\text{ice}=\text{ind}^*(\alpha)$ only
\emph{almost everywhere} for the unique invariant measure --- and a measure-zero set need not be
countable. BHZ's own Proposition 4.1 supplies the counterexample: the ``keep one'' intercept
$c_k=a_k-1$ has $\text{ice}(\omega)\le1+\theta=2.618\ldots$ at \emph{every} irrational slope, while
$\text{ind}^*(\alpha)=2+\limsup_k[a_k;a_{k-1},\dots,a_1]$ is unbounded in $A$; its digit string is
neither eventually $0$ nor either shift-orbit pattern of the characteristic sequence, so its
intercept is not a return intercept. Moreover the family
$\{(c_k)\ \text{admissible}: a_k-c_k\in\{1,2\}\}$ is \emph{uncountable} and has $\text{ice}$
uniformly bounded, so for every $A\ge2$ an uncountable set of intercepts has
$\text{ice}(\omega)<\text{ind}^*(\alpha)$. \emph{Consequence:} intercept $0$ is in fact one of the
\emph{easiest} intercepts ($\text{ice}=\text{ind}^*(\alpha)-1\to\infty$), not the hardest; the
hardest are the keep-one intercepts, where $\text{ice}\to2$ as $A\to\infty$. Revision 1's Open 2 is
withdrawn, its premise being false.

\item[C3 --- Convergent-denominator indexing.] BHZ Proposition 2.7 fixes
$\alpha=[0;a_1+1,a_2,a_3,\dots]$ with $q_0=1$, $q_1=a_1+1$, $q_k=a_kq_{k-1}+q_{k-2}$ for $k\ge2$:
the BHZ \emph{digit} $a_1$ is one less than the first partial quotient of $\alpha$, while $(q_k)$
are the ordinary convergent denominators. This program's record and code used $q_1=a_1$. Every
$\limsup$ is insensitive to a single shifted digit, so \emph{every} $\text{ice}$ value reported in
Revision 1 is correct; every \emph{root length} computed from $q_k$ was not. This is the complete
explanation of the root-length mismatch the Phase 7 record disclosed against itself --- BHZ
Proposition 3.3 was never wrong. Under the correct convention the root-length formulas reproduce
exactly against literal longest common prefixes (562 checks, no failure).

\item[C4 --- Missing length factor in the surplus decomposition.] Revision 1 \S13 displayed
$S_k\ge A\cdot E_k-C_A\cdot D_k-O_A(\log\ell_k)$. The leading term cannot be $A\cdot E_k$: $A$ is a
pure number while $S_k$ scales like a length. Re-derived from the integer periodic-value formula
with every factor kept (\S\ref{sec:ostrowski}), the correct statement on the odd-weight branch is
$S_k\ge\ell_k E_k+1-\lceil D(R_k)\rceil\log_23-\log_2(2\ell_k)$, with $\ell_k=|W_k|$ the
\emph{Sturmian} root length. Revision 1's own flexible condition $E_k\gg\log\ell_k/\ell_k$ is
equivalent to $\ell_kE_k\gg\log\ell_k$ and so is consistent with $\ell_k\cdot E_k$; the displayed
inequality was a transcription slip, not a different claim, and no conclusion drawn from it changes.

\item[C5 --- Conditional attainment of $y(k)$.] Revision 1 treated both Ostrowski expressions as
giving actually-occurring initial powers for every admissible $(c_k)$. BHZ Proposition 3.3 attains
$x(k)$ for \emph{every} $k$, with primitive root of length exactly $q_k$, but attains $y(k)$ only
when $0<c_k<a_k$, with root length $q_k-c_kq_{k-1}$; Corollary 3.5's three observations
($c_k=a_k\Rightarrow y(k)=x(k-2)$; $c_k=0,c_{k+1}=a_{k+1}\Rightarrow y(k)<y(k+1)=x(k-1)$;
$c_k=0,c_{k+1}<a_{k+1}\Rightarrow y(k)\le x(k)$) are what upgrade the $\limsup$. For a $\limsup$
this is harmless; for a proof that must hand the bridge an \emph{attained} root it is not, because
the branch of the argument that is tightest at $A=1$ is exactly $c_k=a_k$, where $y(k)$ is not
directly attained. Theorem \ref{thm:generalmargin} routes every case through an attained witness
explicitly.

\end{description}

\paragraph{What survives unchanged.} Theorem \ref{thm:abstract} (the abstract criterion), the
discrepancy--height bound, the repetition--height competition, the exact XOR transfer
$L_v=L+1$, the density-$1/2$ structural fact, the intercept-$0$ case of
Theorem \ref{thm:boundedtype} and its $\text{term2}(k)>2+1/(A+1)$ analysis, the whole
self-correction history of \S\ref{sec:selfcorrection}, the counting-route material, and every
numerical finding about structured intercepts --- including the identification of the
near-maximal-digit family as the empirical worst case, which \S\ref{sec:generalintercept} now
confirms is \emph{exactly} extremal. What is superseded: Revision 1's \S12 conjecture, its Open 1,
Open 2 and Open 3, and the two displayed formulas of C3 and C4.

\paragraph{Attribution.} The load-bearing symbolic input to the general-intercept theorem is
Berth\'e--Holton--Zamboni's, not this program's. This note claims as its own contribution the
arithmetic application: the periodic-approximation/arithmetic-height bridge, the exact XOR
transfer, the density-$1/2$ fact, the corrected surplus inequality, and the attained-witness
bookkeeping with explicit growing root lengths that turns a $\limsup$ bound on $\text{ice}$ into an
unbounded surplus. Recording a published 2006 theorem as an open conjecture was an error of
literature search, and is corrected here rather than quietly dropped.

\section{The Periodicity Conjecture and the Conjugacy Map}
\label{sec:conjecture}

Let $T:\Z_2\to\Z_2$ be the $3x+1$ map in its standard form,
\[
T(x)=\begin{cases} x/2, & x\equiv 0 \pmod 2,\\ (3x+1)/2, & x\equiv 1\pmod 2.\end{cases}
\]
The \emph{parity vector} of $x\in\Z_2$ is $v(x)=(v_0,v_1,\dots)\in\{0,1\}^\N$ with $v_i\equiv T^i(x)\pmod 2$. Bernstein and Lagarias showed that $x\mapsto v(x)$ is a bijection of $\Z_2$ onto $\{0,1\}^\N$, and in fact a \emph{2-adic isometry} \cite{BernsteinLagarias1996,Bernstein1994}: identifying $\{0,1\}^\N$ with $\Z_2$ through binary expansions,
\[
v_2\bigl(\Phi(a)-\Phi(b)\bigr) = \operatorname{lcp}(a,b),
\]
where $\Phi$ is the inverse map (so $\Phi(v)$ is the unique 2-adic integer with parity vector $v$), and $\operatorname{lcp}$ is longest-common-prefix length. Conjugating $T$ by $\Phi$ turns it into the shift map on $\{0,1\}^\N$. This isometry property is the single fact the entire bridge in this note exploits: two sequences agreeing on a long prefix must have 2-adically close $\Phi$-images.

Rationality of $\Phi(v)$ corresponds to eventual periodicity of $v$ in one direction only: if $v$ is eventually periodic then $\Phi(v)\in\Q$. The converse is:

\begin{quote}
\textbf{Lagarias's Periodicity Conjecture} \cite[\S2]{Lagarias1985}. No aperiodic parity vector has a rational $\Phi$-value.
\end{quote}
(See also López and Stoll's continuous-setting formulation of the same conjecture, \cite{LopezStoll2021}.)

This conjecture is open, and \emph{very little is known unconditionally in the converse direction} --- a fact worth stating plainly, since it frames everything that follows. The conjecture quantifies over \emph{all} $v\in\{0,1\}^\N$ that fail to be eventually periodic: an uncountable family with no assumed combinatorial structure.

\paragraph{Full conjecture versus the classes treated here.} This note --- and the research program it consolidates --- does not address the full conjecture. It proves $\Phi(v)\notin\Q$ for two structurally defined sub-families of aperiodic words (Sturmian words, proved in the companion foundation paper; and complementary symmetric Rote sequences over bounded-type slopes at intercept $0$, proved here), and precisely maps, without resolving, one further natural extension (Rote sequences at arbitrary intercepts). \S\ref{sec:open} records a structural fact showing that no extension of this program's method, however far pushed, can by itself close the full conjecture. Nothing in this note should be read as narrowing that gap beyond what is explicitly proved.

\section{Two Routes to Irrationality}
\label{sec:tworoutes}

The foundation paper and this program's extension of it use two logically independent mechanisms to force $\Phi(v)\notin\Q$.

\subsection{Route A --- Counting}
\label{sec:routeA}

Let $p_v(L)$ denote the number of distinct length-$L$ factors of $v$ (its \emph{factor complexity}). Dubickas proved in 2009 \cite{Dubickas2009} that a positive-integer $3x+1$ trajectory tending to infinity has parity-word complexity $p(L) > 1.70951129\,L$ eventually; the foundation paper transports this to every rational 2-adic integer, since $T$ preserves a fixed odd denominator. The exact threshold constant is
\[
\frac{1}{\log_2(3/2)} = 1.709511\ldots,
\]
and the transported statement is: if $v$ is not eventually periodic and $\Phi(v)\in\Q$, then $\liminf_L p_v(L)/L \ge 1/\log_2(3/2)$. Contrapositively: \emph{any} aperiodic $v$ with $\liminf_L p_v(L)/L < 1.70951\ldots$ has $\Phi(v)\notin\Q$, unconditionally, with no repetition or discrepancy hypothesis whatsoever.

This is a purely counting-based exclusion, and it is exactly as strong as its complexity bound: it says nothing once a word's complexity crosses the threshold. Sturmian words ($p(n)=n+1$) sit far below it. Complementary symmetric Rote sequences, the next class this program treats, have $p(n)=2n$ --- strictly above $1.70951\ldots\,n$ for every $n\ge 2$ --- and are therefore \emph{entirely outside Route A's reach}. This is precisely where Route A stops, and why a second, structurally different mechanism is needed to make any further progress.

\subsection{Route B --- Periodic Approximation and Arithmetic Height}
\label{sec:routeB}

This is the central mechanism of this note. Informally, the chain of implications is: symbolic periodic approximation $\to$ exact 2-adic agreement $\to$ divisibility $\to$ rational periodic approximant $\to$ Archimedean height $\to$ contradiction.

Concretely: if $v$ agrees with a periodic word $W^\infty$ on a long prefix, the 2-adic isometry forces $\Phi(v)$ and the \emph{rational number} $\Phi(W^\infty)$ to be 2-adically close to a precisely quantified depth. If $\Phi(v)$ were itself rational, this closeness becomes an exact integer divisibility statement between two specific integers built from $\Phi(v)$'s numerator/denominator and $W$'s combinatorial data. When the periodic approximation is long and clean enough relative to the arithmetic size (\emph{height}) of the numbers involved, that divisibility is impossible unless the two rationals coincide --- and they were assumed not to. The whole of \S\ref{sec:abstract}--\S\ref{sec:repheight} makes this precise.

\section{The Abstract Bridge Theorem}
\label{sec:abstract}

Following Proposition 4.1 of the foundation paper, every finite word $W$ (length $\ell=|W|$, $k=|W|_1$ ones) has an explicit periodic-value formula
\[
\Phi(W^\infty) = \frac{c_W}{2^\ell - 3^k}
\]
for an explicit integer $c_W$ (\repo{THEOREM\_STATUS.md}, citing the foundation paper's Prop.\ 4.1).

Suppose, for contradiction, $\Phi(s)=u/v\in\Q\cap\Z_2$ in lowest terms. For any finite word $W_n$ with $s\ne W_n^\infty$, put $\ell_n=|W_n|$, $k_n=|W_n|_1$, $L_n=\operatorname{lcp}(s,W_n^\infty)$, $\delta_n = 2^{\ell_n}-3^{k_n}$, and define the integer
\[
M_n := u\,\delta_n - v\,c_{W_n}.
\]
By the 2-adic isometry, $v_2(\Phi(s)-\Phi(W_n^\infty))\ge L_n$, which translates into a 2-adic divisibility lower bound on $M_n$; on the other hand $|M_n|$ is bounded above in terms of $|u|,v,\delta_n,c_{W_n}$ --- an Archimedean (ordinary absolute-value) bound. When the 2-adic lower bound forces $M_n$ to be divisible by more than its Archimedean bound allows, $M_n=0$ is forced, and (with $s\ne W_n^\infty$) that is eventually impossible. This comparison is exactly what the repository's \emph{surplus} quantity measures.

\begin{definition}[Surplus, canonical form \cite{Phase2Surplus}]
\[
S_s(W) := \operatorname{lcp}(s,W^\infty) - \log_2\bigl(|2^{|W|}-3^{|W|_1}| + c_W\bigr).
\]
\end{definition}
This is the repository's own proof-native quantity, fixed from Phase 1 onward and not altered by any later phase: \emph{``the proof-native quantity, not a convenient after-the-fact packaging''} (\repo{theorems/phase2/SURPLUS\_THEOREM.md}).

\begin{theorem}[Abstract 2-adic periodic-approximation criterion \cite{Phase2Abstract}]
\label{thm:abstract}
Let $s\in\{0,1\}^\N$ and $(W_n)_{n\ge1}$ finite binary words with $s\ne W_n^\infty$ for all $n$. If
\[
\limsup_{n\to\infty} S_s(W_n) = +\infty,
\]
then $\Phi(s)\notin\Q$.
\end{theorem}
This is stated and proved with no Sturmian or Rote terminology at all --- it is the Sturmian-free abstraction every later result in this note specializes.

\begin{theorem}[Dual necessity, \cite{Phase2Surplus}]
If $\Phi(s)=u/v\in\Q$ in lowest terms, height $H=\max(|u|,v)$, then for every finite $W$ with $s\ne W^\infty$,
\[
S_s(W) \le \log_2 H.
\]
\end{theorem}
Together: \emph{surplus $\to+\infty$ along some sequence of periodic approximants forces irrationality; rationality forces surplus to stay uniformly bounded.} Everything that follows is the construction of periodic approximants along which surplus is provably unbounded, for successively broader symbolic classes.

\section{Discrepancy as the Height Bridge}
\label{sec:discrepancy}

Repetition length $L_n$ alone is not enough: the surplus definition also involves $c_W$, and the periodic-value formula gives no a priori control on how large $c_W$ can be relative to $2^\ell-3^k$. The quantity that closes this gap is prefix \emph{discrepancy}:
\[
D(W) := \max_{0\le i\le \ell} \left| k_i(W) - \frac{ik}{\ell}\right|,
\]
where $k_i(W)$ counts ones in the length-$i$ prefix of $W$.

\begin{theorem}[Discrepancy-height bound \cite{Phase1Discrepancy}]
For every finite word $W$, $\ell=|W|$, $k$ ones, $D=D(W)$:
\[
c_W \;\le\; \ell\cdot 3^{\lceil D\rceil}\cdot\max(2^\ell,3^k).
\]
\end{theorem}
(Proved, and exhaustively verified by direct computation for $\ell\le 16$.) This single inequality is what converts symbolic regularity (a bound on how unevenly $1$s are distributed in $W$) into an arithmetic height estimate the surplus computation can use. The relevant hierarchy of discrepancy regimes, from weakest to strongest sufficient hypothesis actually used somewhere in this program, is:
\begin{itemize}
\item $D(W)=O(1)$ --- e.g., single-block or highly regular words; strongest, rarely needed.
\item $D(W)=O(\log \ell)$ --- achieved for rotations by any slope with \emph{bounded partial quotients}, via the classical three-distance theorem; this is the regime used throughout \S\ref{sec:oddweight}--\S\ref{sec:boundedtype}.
\item $D(W)=o(\ell)$ --- the general sufficient hypothesis for the abstract sublinear-discrepancy theorem below; strictly weaker, but gives up the explicit logarithmic rate.
\end{itemize}
Sublinear discrepancy is exactly what is needed for the height cost to be dominated by a linear repetition gain (\S\ref{sec:repheight}); the sharper $O(\log\ell)$ rate is what the bounded-partial-quotient hypothesis buys on top of that.

\section{Repetition--Height Competition}
\label{sec:repheight}

\begin{theorem}[Sublinear-discrepancy repetition criterion \cite{Phase1Discrepancy}]
Suppose $s$ begins in initial powers $W_n^{r_n}$, $\ell_n=|W_n|\to\infty$, densities $\gamma_n=k_n/\ell_n\to\gamma_\infty$, and $D(W_n)=o(\ell_n)$. If
\[
\liminf_n r_n > \max(1,\gamma_\infty\log_2 3) \quad\text{strictly},
\]
then $S_s(W_n)\to+\infty$, hence $\Phi(s)\notin\Q$.
\end{theorem}
Taking the supremum over $\gamma\in(0,1)$ of $\max(1,\gamma\log_2 3)$ gives $\log_2 3=1.58496\ldots$, so $r>\log_2 3$ suffices \emph{uniformly in the density}. Because $\log_2 3<2$, an initial \emph{square} ($r=2$) always clears this bar, with an explicit margin $2-\log_2 3=0.41503\ldots$ --- which is exactly why squares are the workhorse repetition structure throughout this program:

\begin{theorem}[Squares suffice unconditionally \cite{Phase1Repetition}]
If $s$ begins in arbitrarily long initial squares $W_nW_n$ --- with \emph{no} hypothesis on discrepancy, balance, or internal structure of the roots $W_n$ --- then $\Phi(s)\notin\Q$.
\end{theorem}

The later, more flexible sufficient condition used once a fixed margin is no longer available (\S\ref{sec:ostrowski}) does not require a fixed constant margin at all; it only requires the repetition \emph{excess} above the threshold to dominate the discrepancy-plus-logarithmic height cost, i.e.\ excess $\cdot\, \ell \gg \log \ell$ rather than excess $\ge$ const.

\section{The Sturmian Foundation}
\label{sec:sturmian}

The foundation paper \cite{DeJesus2026Sturmian} proves, by \emph{two independent mechanisms}, that every Sturmian parity word has irrational $\Phi$-value:

\begin{enumerate}
\item \textbf{Counting.} Every Sturmian word has $p(L)=L+1$, far below the Dubickas threshold, so Route A (\S\ref{sec:routeA}) applies directly.
\item \textbf{Liouville / initial-square route.} Sturmian words have arbitrarily long initial squares (Theorem 10.2 of \cite{DeJesus2026Sturmian}, citing Allouche, Davison, Queffélec, and Zamboni \cite{ADQZ2001}), which triggers the repetition--height mechanism of \S\ref{sec:repheight} directly (Theorem 10.1 of \cite{DeJesus2026Sturmian}). The foundation paper also gives a second, independent and effective proof, for the specific case López and Stoll's own paper \cite{LopezStoll2009} had identified as open (their critical-density Sturmian word at slope $\beta=\ln2/\ln3$): this is then generalized to every irrational slope, and a further argument removes the restriction to intercept $0$ via an initial square of the word itself.
\end{enumerate}

Both mechanisms are cited, not re-derived, in this note; nothing in this program alters or extends the Sturmian case itself. What matters for everything that follows is where the foundation paper's own abstract explicitly places the frontier: \emph{``Counting reaches every complexity below $1.70951\ldots L$; the frontier is complexity $2L$ (Rote words, rotation codings), open for both routes.''} This is the exact motivation for the program consolidated here.

\section{Crossing the Counting Frontier: Complementary Symmetric Rote Sequences}
\label{sec:rote}

A \emph{complementary symmetric (CS) Rote sequence} $v\in\{0,1\}^\N$ is one whose first-difference (XOR) transform is Sturmian:
\[
v \text{ is CS Rote} \iff u:=S(v) \text{ is Sturmian}, \qquad u_i = v_i\oplus v_{i+1}\ \ (i\ge0),
\]
with inverse recovery $v_{n+1}=v_n\oplus u_n$ \cite{Phase2Transfer}. Such sequences have factor complexity exactly $p(n)=2n$, standard in the combinatorics-on-words literature on this class (background on the surrounding Sturmian/complementary-word theory: \cite{Lothaire2002}; see the bibliography note on primary Rote-sequence sourcing) --- past the Dubickas threshold, so Route A never applies to them; every irrationality result about CS Rote sequences in this note is a Route B result.

\begin{theorem}[Exact transfer, with a $+1$ bonus \cite{Phase2Transfer}]
\label{thm:transfer}
Let $L:=\operatorname{lcp}(u,W^\infty)$ for a Sturmian $u=S(v)$ and finite word $W$, and let $L_v:=\operatorname{lcp}(v,\mathrm{target}^\infty)$ where $\mathrm{target}=V$ (if $W=V$ has even weight) or $\mathrm{target}=R:=V\bar V$ (if $V$ has odd weight, $\bar V$ its bitwise complement). Then
\[
L_v = L+1, \quad\text{exactly, not merely } \ge L \text{ or } \approx L.
\]
\end{theorem}
Verified by induction (both directions of the inequality) and computationally over $6$ initial cases in Phase 2; re-derived completely independently in Phase 5 and cross-checked exhaustively over every binary root of length $2$ to $10$ and every mismatch position from $\ell$ to $4\ell$ --- $114{,}680$ cases, $0$ failures \cite{Phase5Transfer}. The theorem is intercept-agnostic: it holds verbatim regardless of the intercept of the underlying Sturmian word.

\section{The Odd-Weight Mechanism}
\label{sec:oddweight}

This is the structural discovery that makes the whole bounded-type theorem (\S\ref{sec:boundedtype}) possible, and it deserves to be isolated on its own.

When the Sturmian root $W$ has \emph{odd} weight, Theorem \ref{thm:transfer}'s target is $R=V\bar V$, not $V$ itself.

\begin{theorem}[Density of the doubled root, exactly $1/2$ \cite{Phase4Required}]
For every finite word $V$, $R:=V\bar V$ satisfies $|R|=2|V|$ and $|R|_1=|V|$, hence
\[
\text{density}(R) = \frac{|R|_1}{|R|} = \frac12, \quad\text{exactly, unconditionally.}
\]
\end{theorem}
\begin{proof}[Proof (as recorded)]
$|R|_1 = |V|_1 + |\bar V|_1 = |V|_1 + (|V|-|V|_1) = |V| = \ell$; and $|R|=2\ell$; so density $=\ell/(2\ell)=1/2$. $\blacksquare$
\end{proof}

Arithmetically, this matters because $3^{|R|_1}=3^{|R|/2}$ grows strictly more slowly than $2^{|R|}$ (since $\sqrt3<2$), so the dominant exponential height scale in the periodic-value denominator $2^{|R|}-3^{|R|_1}$ is unambiguously $2^{|R|}$, regardless of the fine structure of $V$. Plugging $\gamma=1/2$ into the repetition--height criterion of \S\ref{sec:repheight} collapses the required threshold on the Sturmian side from an earlier, over-conservative estimate of $2\log_2 3\approx3.170$ down to the true requirement,
\[
r_{\text{required}}(\gamma) = 2, \qquad\text{uniformly, for both the even- and odd-weight transfer routes.}
\]
This is established in \cite{Phase4Required}. This single fact replaced two working hypotheses this program had earlier formed and found false --- a universal even-weight witness hypothesis, and a silver-ratio extremality claim --- both recorded in full in \S\ref{sec:selfcorrection}.

\section{Self-Correction History}
\label{sec:selfcorrection}

This program is falsification-first throughout: every phase actively attacked its own prior claims before extending them. Four corrections are recorded here in full, not hidden, because they show how the final theorem (\S\ref{sec:boundedtype}) was actually reached, and because a future researcher extending this work needs to know which paths were tried and ruled out.

\paragraph{A. Universal even-weight witness hypothesis --- rejected.} An early working assumption was that the Rote-transfer route could always be arranged to use an even-weight root $V$ (avoiding the doubled root $R$ altogether). A finite $8$-state automaton analysis of the weight-parity sequence found exactly one ``bad'' cycle sustaining all-odd weight forever, and an explicit quadratic-irrational counterexample slope was found and verified to $28$ convergents: $\gamma=[0;1,1,1,\overline{2}]$ has weight parity $0,0,1,1,1,1,1,\ldots$ --- all odd (weight $1$) from index $2$ onward, forever \cite{Phase3Automaton}. The universal claim is false. The bridge still succeeds on this slope because the odd-weight transfer route (\S\ref{sec:oddweight}) handles exactly this case.

\paragraph{B. Silver-ratio extremality --- rejected, explicit counterexample.} A subsequent working hypothesis, formed from testing only four continued-fraction tails, was that among all-even-tail slopes the silver ratio ($\overline{2}$ tail) is \emph{extremal}, with minimal achievable exponent $2+\sqrt2\approx3.4142$. An adversarial search found $\gamma=[0;1,4,1,5,1,\overline{6,2}]$ with $\text{ice}(u)\approx3.1547 < 2+\sqrt2$ \cite{Phase4Silver}, disproving the hypothesis outright. The constant $2+\sqrt2$ itself is not wrong --- it is the exact value for the pure silver-ratio tail \cite{Phase4SilverAudit} --- but it was the wrong threshold to be solving for: once the density-$1/2$ fact (\S\ref{sec:oddweight}) is known, the actual required threshold is $r>2$, and $2+\sqrt2\gg2$ was never close to the real boundary.

\paragraph{C. Seed-independence via complement-invariance of $c_W$ --- rejected, repaired.} The claim that both CS Rote seeds ($v_0=0$ or $v_0=1$) give identical surplus was correct, but its stated justification (``complementing uniformly complements $R$, so $c_W$ is unaffected'') was not: direct testing found $c_W\ne c_{\bar W}$ in $8$ of $8$ random odd-weight cases, often by an order of magnitude \cite{Phase5Hostile}. The repaired reason: every verification script's height bound $F_{\text{bound}}(W)$ depends on the root only through $(\ell,k,D(\text{root}))$, never through $c_W$ directly, and $D(V\bar V)=D(\bar V V)$ exactly, checked in $15$ of $15$ random cases. The \emph{conclusion} survives untouched; only its justification needed repair --- exactly the kind of error a hostile self-audit is designed to catch.

\paragraph{D. Phase 6 small-$k$ boundary correction.} The load-bearing inequality $\text{term2}(k)>2+1/(A+1)$ (\S\ref{sec:boundedtype}) was stated to hold strictly for every $k$. Adversarial testing against genuinely non-eventually-periodic bounded-type controls (Thue--Morse-coded, Fibonacci-word-coded, and two independent random $\{1,2\}$ sequences) found one case of non-strict \emph{equality} rather than a violation: in the Fibonacci-coded control, $\text{term2}(2)$ equals $2+1/(A+1)=7/3$ exactly, at $k=2$, the only such case out of $232$ checks across all four sequences \cite{Phase6Arrow}. This does not weaken the theorem, because the abstract criterion (Theorem \ref{thm:abstract}) only requires $\limsup$ unboundedness along \emph{some} growing sequence of witnesses, not strict inequality at literally every finite $k$.

None of these four corrections changed the truth of the final headline theorem; each changed (and, in cases B and C, simplified) the reason it is true.

\section{The Bounded-Type CS Rote Theorem (Every Intercept)}
\label{sec:boundedtype}

\begin{theorem}[Headline result of the program \cite{Phase8Proof}]
\label{thm:boundedtype}
Let $\gamma\in(0,1)$ be irrational with bounded partial quotients ($A:=\sup_n a_n<\infty$). Let $u$
be \textbf{any} Sturmian word of slope $\gamma$ --- \textbf{any} intercept $\rho$, either the upper
or the lower mechanical convention. Let $v$ be either complementary symmetric Rote sequence with
$S(v)=u$ (either seed), or any shift $\sigma^j(v)$. Then
\[
\Phi(v)\notin\Q.
\]
\end{theorem}

The intercept-$0$ case of Theorem \ref{thm:boundedtype} is the result Revision 1 stated
\cite{Phase6Bounded}; the general-intercept extension is \S\ref{sec:generalintercept}, whose
symbolic input is BHZ \cite{BHZ2006} \S5 (Revision 2, C1). The proof architecture of
\S\ref{sec:architecture} is unchanged apart from that one arrow: every other step was already
proved intercept-agnostic, and each has been re-verified on literally constructed mechanical words
at general intercepts, under both conventions and both seeds \cite{Phase8Chain}.

This supersedes an earlier theorem of the same program, proved first for the strictly narrower class of \emph{quadratic}-irrational slopes \cite{Phase4Quadratic}, \emph{without weakening any of its quantifiers} --- both seeds, every shift, still covered. The upgrade matters because quadratic irrationals are countable (each is a root of an integer quadratic, and there are only countably many such polynomials), while bounded-partial-quotient irrationals are uncountable: for each bound $A$, the set of irrationals whose partial quotients never exceed $A$ contains a Cantor-type subset, and the union over all $A$ is an uncountable set.

\begin{remark}[What ``uncountable'' does and does not say here]
Uncountability is a statement about cardinality only. The set of bounded-type irrationals has Lebesgue measure zero: a Lebesgue-almost-every real number has \emph{unbounded} partial quotients. Theorem \ref{thm:boundedtype} says nothing about a randomly chosen real slope, and nothing about density among all irrationals --- only that the class it covers, while structurally constrained, is not merely a countable list of algebraic special cases.
\end{remark}

\section{Proof Architecture of the Bounded-Type Theorem}
\label{sec:architecture}

The shortest audited chain, reconstructed in Phase 6 without any periodicity hypothesis and cross-checked against four genuinely non-eventually-periodic bounded-type controls \cite{Phase6Arrow}, adapting the original Phase 4 chain \cite{Phase4Chain}:

\begin{center}
\begin{tabular}{ll}
\toprule
Step & Cites \\
\midrule
bounded partial quotients, $A=\sup_n a_n$ & hypothesis \\
$\downarrow$ Sturmian words have arbitrarily long initial squares & Thm 10.2 \cite{DeJesus2026Sturmian}, \cite{ADQZ2001} \\
$\downarrow$ BHZ margin at intercept $0$, $\text{term2}(k)>2+\tfrac1{A+1}$ (every $k$, up to \S\ref{sec:selfcorrection}.D) & \cite{BHZ2006}, Cor.\ 3.5 \\
$\downarrow$ \emph{or, at any intercept:} $\text{ice}(\omega)\ge2+\delta(A)$, attained, roots $\to\infty$ & \cite{BHZ2006}, Prop.\ 5.1--5.2; Thm \ref{thm:generalmargin} \\
$\downarrow$ arbitrarily long usable witnesses & \S\ref{sec:repheight} \\
$\downarrow$ exact XOR transfer, $L_v=L+1$ & Thm \ref{thm:transfer} \cite{Phase2Transfer,Phase5Transfer} \\
$\downarrow$ even/odd branch split & \S\ref{sec:rote} \\
$\downarrow$ odd branch: doubled root density exactly $1/2$ & \S\ref{sec:oddweight} \cite{Phase4Required} \\
$\downarrow$ actual-root discrepancy control, $O(\log\ell)$, uniform in $\rho$ & \cite{Phase2Discrepancy}, repaired \cite{Phase5Hostile}, Lem.\ \ref{lem:halfslope} \\
$\downarrow$ sublinear height penalty & \S\ref{sec:discrepancy} \\
$\downarrow$ positive, unbounded surplus & \S\ref{sec:abstract} \\
$\downarrow$ abstract irrationality criterion (Thm \ref{thm:abstract}) & \cite{Phase2Abstract} \\
$\Rightarrow$ $\Phi(v)\notin\Q$ & Thm \ref{thm:boundedtype} \\
\bottomrule
\end{tabular}
\end{center}

Notably absent from this chain: the weight-parity automaton of \S\ref{sec:selfcorrection}.A, and the silver-ratio constant $2+\sqrt2$ of \S\ref{sec:selfcorrection}.B --- both were real, correctly-derived facts about this problem, and neither is needed by the final proof.

\section{General Intercepts}
\label{sec:generalintercept}

Sections \ref{sec:oddweight}--\ref{sec:architecture} were originally stated for the
\emph{intercept-$0$} representative of a slope $\gamma$: the lower mechanical word
$u_i=\lfloor(i+1)\gamma\rfloor-\lfloor i\gamma\rfloor$. This section removes that restriction.

An arbitrary intercept is encoded by an \emph{Ostrowski digit sequence} $(c_k)_{k\ge1}$, subject to
BHZ's admissibility condition \cite{BHZ2006}:
\[
a_k \ge c_k \ge 0 \text{ for all } k,\quad a_k\ge1 \text{ for } k\ge2,\quad \text{and if } c_k=a_k \text{ then } c_{k-1}=0.
\]
The characteristic (intercept-$0$) sequence is $c_k\equiv0$. Where the partial quotients $(a_k)$ fix
the canonical Sturmian \emph{scale hierarchy} (the convergent denominators $q_k$), the digits
$(c_k)$ determine \emph{where inside that hierarchy} a particular intercept's word sits.

\paragraph{Conventions, stated exactly (Revision 2, C3).} BHZ Proposition 2.7 sets
\[
\alpha=[0;a_1+1,a_2,a_3,\dots],\qquad q_0=1,\quad q_1=a_1+1,\quad q_k=a_kq_{k-1}+q_{k-2}\ (k\ge2),
\]
so the BHZ digit $a_1$ is one less than the first partial quotient of $\alpha$, while $(q_k)$ are
the ordinary convergent denominators. Write $A_b:=\sup_k a_k$; always $A_b\le A$. Bounds below are
sharpest in $A_b$, and remain true with $A$ substituted.

At intercept $0$, BHZ's Corollary 3.5 collapses to the clean closed form used in
\S\ref{sec:architecture}. In general it does not collapse:
\[
x'(k) = \frac{\sum_{j=1}^{k+1}(a_j-c_j)q_{j-1}}{q_k}, \qquad
y(k) = 1 + \frac{\sum_{j=1}^{k}(a_j-c_j)q_{j-1}}{q_k-c_kq_{k-1}},
\]
\[
\text{ice}(\omega)=\limsup_k \max(x'(k),y(k)).
\]
Both terms depend on the entire digit history, not just on $(a_j)$.

\paragraph{Attainment, stated exactly (Revision 2, C5).} By BHZ Proposition 3.3, $x(k)$ is the
prefix power of an actually occurring initial power of $\omega$ with primitive root of length
exactly $q_k$, \emph{for every} $k$; and $y(k)$ is, with root length exactly $q_k-c_kq_{k-1}$,
\emph{provided} $0<c_k<a_k$. Nothing is asserted about $y(k)$ as an attained power otherwise;
Corollary 3.5's three observations convert those cases into attained $x$-witnesses.

\paragraph{The margin is a published theorem (Revision 2, C1).}

\begin{theorem}[Berth\'e--Holton--Zamboni \cite{BHZ2006}, Propositions 5.1 and 5.2]
\label{thm:bhzmargin}
If $(s,t)$ are integers with $s>1$ and $(a_k,a_{k+1})=(s,t)$ for infinitely many $k$, then every
$\omega\in X_\alpha$ satisfies $\text{ice}(\omega)\ge2+\frac1{2(s+1)(t+1)+1}$. If
$(a_k,a_{k+1},a_{k+2})=(1,1,t)$ for infinitely many $k$, then every $\omega\in X_\alpha$ satisfies
$\text{ice}(\omega)\ge2+\frac1{8t+1}$.
\end{theorem}

\begin{corollary}
\label{cor:bhzbounded}
For every irrational $\gamma$ of bounded type and every $\omega\in X_\gamma$,
$\text{ice}(\omega)\ge2+\frac1{2(A_b+1)^2+1}$.
\end{corollary}

\begin{proof}
There are finitely many digit pairs. If $a_k>1$ for infinitely many $k$, some pair $(s,t)$ with
$s>1$ and $s,t\le A_b$ recurs infinitely often, and Proposition 5.1 applies. Otherwise $a_k=1$ for
all large $k$, so $(a_k,a_{k+1},a_{k+2})=(1,1,1)$ recurs infinitely often and Proposition 5.2
applies with $t=1$, giving $2+\frac19\ge2+\frac1{2(A_b+1)^2+1}$.
\end{proof}

Here $X_\gamma$ is BHZ's set of \emph{all} Sturmian sequences of slope $\gamma$, i.e.\ the
itineraries of every point under either of the two interval exchanges of their \S2.1 --- so
Corollary \ref{cor:bhzbounded} quantifies over every intercept \emph{and} both mechanical
conventions at once, which is exactly the quantifier the bridge needs.

\paragraph{An attained-witness form, with growing roots.} Corollary \ref{cor:bhzbounded} bounds a
$\limsup$. A $\limsup$ bound $\text{ice}(\omega)\ge2+\delta$ yields, for each $\varepsilon>0$,
infinitely many $n$ with an attained initial power of exponent $>2+\delta-\varepsilon$; taking
$\varepsilon=\delta/2$ already supplies attained witnesses of exponent $>2+\delta/2$ whose root
lengths necessarily tend to infinity, which suffices for Theorem \ref{thm:boundedtype}. The
following stronger, per-scale form is what this program proves independently, and is what makes the
surplus bound uniform in $k$ rather than only along an unspecified subsequence.

\begin{theorem}[Attained general-intercept margin \cite{Phase8Proof}]
\label{thm:generalmargin}
Let $\gamma$ be irrational with $A_b<\infty$ and let $\omega$ be any Sturmian sequence of slope
$\gamma$ (any intercept, either convention). Put $\delta(A_b):=\frac1{2(A_b+1)^2}$. Then for every
$k\ge7$ there is a word $W$ with $|W|\ge q_{k-4}$ such that $W^r$ is a prefix of $\omega$ with
$r\ge2+\delta(A_b)$. In particular $\text{ice}(\omega)\ge2+\delta(A_b)$, and the bound is realised by
actually occurring initial powers whose root lengths tend to infinity.
\end{theorem}

The proof is a three-case analysis on $d_k:=a_k-c_k$ driven by the substitution
$\lambda_k:=q_{k-1}/q_k$, $m_k:=t_{k-1}+d_k$ with $t_k:=\lambda_km_k$, under which
$x'(k)=m_{k+1}$, $q_k-c_kq_{k-1}=q_{k-1}(\lambda_{k-1}+d_k)$ and
\[
y(k)-2=\frac{\lambda_{k-1}\,(m_{k-1}-1)}{\lambda_{k-1}+d_k},
\]
so that $y(k)>2$ if and only if $m_{k-1}>1$. Admissibility forbids two consecutive $d_k=0$, which
floors $m_k$; the three cases $d_k=0$ (route through $x(k-2)=y(k)$), $d_k\ge2$ (route through
$x(k-1)\ge m_k$) and $d_k=1$ (route through $y(k)$, or through $x(k)$/$x(k-1)$ when $c_k=0$) each
deliver an \emph{attained} power. Full proof in \cite{Phase8Proof}; BHZ's own three-case proof of
Proposition 5.1 is structurally the same trichotomy, reached independently, with constants agreeing
to a factor $1+\frac1{2(A_b+1)^2}$.

\paragraph{How sharp is the margin.} Let $\beta_A$ denote the unique root $>1$ of
$\beta^2=A\beta+1$, i.e.\ $\beta_A=\frac{A+\sqrt{A^2+4}}2$; equivalently
$\beta_A=\lim_kq_k/q_{k-1}$ for any slope whose partial quotients are eventually $A$. On BHZ's
``keep one'' intercept $c_k=a_k-1$ (so $d_k\equiv1$) with BHZ digits eventually $\equiv A$, the
recurrences above solve in closed form: $\lambda_k\to1/\beta_A$, $m_k\to\beta_A/(\beta_A-1)$, and
since $\beta_A^2-1=A\beta_A$,
\[
\text{ice}(\omega)=2+\frac1{A\beta_A}\qquad\text{exactly.}
\]
Hence $\delta^*(\gamma)\le\frac1{A\beta_A}=\Theta(A^{-2})$ there, so the proved floor
$\frac1{2(A_b+1)^2}$ is off by at most a bounded factor (between $2$ and $4.90$), and --- because
$1/(A\beta_A)\to0$ --- \textbf{no margin uniform in $A$ can exist}: the $A$-dependence in
Revision 1's Conjecture 12.1 was necessary, not a convenience. This also recovers BHZ Proposition
4.1 from the other side, their bound $\text{ice}(\text{keep one})\le1+\theta$ and the floor above
pinning the keep-one value between $2+\frac1{A\beta_A}$ and $1+\theta$.

\begin{remark}[What is \emph{not} claimed]
That $\frac1{A\beta_A}$ \emph{is} the infimum of $\text{ice}-2$ over all intercepts is proved only
for slopes whose BHZ digits are eventually constant with $A\ge2$ (a three-case argument: deviating
from $d_k\equiv1$ costs at least $\min(1+A,\,2+1/\beta_A)$, both strictly larger), plus the
Fibonacci slope via BHZ Proposition 4.3. On every other bounded-type slope it is numerical evidence
only and is deliberately not promoted to a theorem \cite{Phase8Margin}. Theorem
\ref{thm:boundedtype} uses none of it: any fixed $\delta>0$ suffices, and Corollary
\ref{cor:bhzbounded} alone would do.
\end{remark}

\paragraph{The remaining intercept-dependent step, and its hidden hypothesis.} Every ingredient of
\S\ref{sec:architecture} after ``a witness exists'' was already proved intercept-agnostic: the XOR
transfer uses no property of $u$ beyond being a binary sequence; the doubled root's density is
$1/2$ by construction; and the discrepancy bound is a statement about a rotation's equidistribution
rate, uniform in the starting phase. That last step needs one hypothesis Revision 1 left implicit.
For the intercept-$\rho$ word, $k_i(u)=\lfloor i\gamma+\rho\rfloor-\lfloor\rho\rfloor$, so
$k_i\bmod2$ codes the orbit $\{i(\gamma/2)+\rho/2\}$ against $[\tfrac12,1)$: the relevant rotation
is by $\gamma/2$, not $\gamma$.

\begin{lemma}
\label{lem:halfslope}
If $\gamma$ has bounded partial quotients then so does $\gamma/2$, with a bound depending only on
$A$.
\end{lemma}

\begin{proof}
Bounded type is equivalent to $c:=\inf_{q\ge1}q\|q\gamma\|>0$, and $c\ge1/(A+2)$. For any $q\ge1$,
$\|q\gamma\|=\|2\cdot q(\gamma/2)\|\le2\|q(\gamma/2)\|$, so
$q\|q(\gamma/2)\|\ge q\|q\gamma\|/2\ge c/2>0$. Hence $\gamma/2$ is badly approximable with constant
at least $1/(2(A+2))$, and its partial quotients are bounded by $\lfloor2(A+2)\rfloor+1$.
\end{proof}

With Lemma \ref{lem:halfslope} the three-distance bound applies to $\gamma/2$ and gives
$D(R_k)=O_A(\log\ell_k)$ uniformly in $\rho$, completing the chain. Only $D=o(\ell)$ is actually
needed; measured across 534 witnesses at general intercepts, $D(R)/\log_2\ell_k$ stays in
$[0.30,0.61]$ while $D(R)/\ell_k$ falls from $0.20$ to $0.0004$ \cite{Phase8Chain}.

\paragraph{Independent verification.} The whole chain was re-run on literally constructed mechanical
words --- certified rational brackets for both $\gamma$ and $\rho$, every floor and ceiling either
decided by the bracket or raised, and no BHZ formula used anywhere --- across four slopes
(including a non-eventually-periodic one), four intercept families (including keep-one), both
mechanical conventions and both seeds: 534 attained witnesses with $\ell$ from $8$ to $10\,946$;
the transfer identity $L_v=L_u+1$ held $534/534$; $2^{L_v}>F(R)$ as exact integers $534/534$; and
the odd-branch density was exactly $1/2$ in all $338$ odd-branch cases \cite{Phase8Chain}.

\section{Ostrowski Repetition Structure}
\label{sec:ostrowski}

A subtlety worth making explicit for anyone continuing this work: \textbf{low factor complexity is not the same thing as strong initial repetition.} A CS Rote sequence may remain globally $p(n)=2n$ while its intercept changes \emph{which} large canonical blocks (Ostrowski/continuant-scale factors) happen to align at the origin --- which is exactly what determines how long an initial repetition is actually available to feed the bridge.

The scale-dependent quantity that tracks this is an \emph{excess}:
\[
E_k := \max(r'_k, r_k) - 2,
\]
where $r'_k,r_k$ are the two branch-specific repetition exponents at scale $k$ \cite{Phase7Surplus}. The corresponding surplus decomposition, re-derived from the integer periodic-value formula with
every length factor kept explicit (Revision 2, C4), is
\[
\text{odd branch:}\quad S_k \;\ge\; \ell_k\,E_k \;+\;1\;-\;\lceil D(R_k)\rceil\log_23\;-\;\log_2(2\ell_k),
\]
\[
\text{even branch:}\quad S_k \;\ge\; \ell_k\,(r_k-\log_23)\;+\;1\;-\;\lceil D(V_k)\rceil\log_23\;-\;\log_2\ell_k,
\]
where $\ell_k=|W_k|$ is the \emph{Sturmian} root length \cite{Phase8Proof}. Revision 1 displayed
$S_k\ge A\cdot E_k-C_A D_k-O_A(\log\ell_k)$, which is dimensionally incomplete: the leading term is
a length times an excess, not the partial-quotient bound times an excess. The flexible sufficient
condition this leads to --- used in place of a fixed margin whenever one is not cleanly available ---
is
\[
E_k \gg \frac{\log \ell_k}{\ell_k}, \qquad\text{equivalently}\qquad \ell_kE_k\gg\log\ell_k,
\]
which is manifestly the same statement as the corrected inequality, and is why no conclusion Revision
1 drew from the displayed formula changes \cite{Phase7Flexible}. With Theorem
\ref{thm:generalmargin} the excess is bounded below by the \emph{constant} $\delta(A_b)$ and
$\ell_k\ge q_{k-4}\to\infty$, so the flexible machinery is not needed at all: the general-intercept
case lands in the easiest regime, $E_k\asymp1$. Every structured intercept family tested lands there
too \cite{Phase7Structured}, and the near-maximal family those tests identified as worst is exactly
the keep-one family, for which the excess is exactly $1/(A\beta_A)$.

The even branch needs no discrepancy hypothesis at all: the combinatorial cap
$D(V)\le\gamma_V(1-\gamma_V)\ell$ together with
$\sup_\gamma[\max(1,\gamma\log_23)+\gamma(1-\gamma)\log_23]=\log_23<2$ already suffices
\cite{Phase1Repetition}. Only the odd branch needs $D=o(\ell)$.

\section{Interpretation: Symbolic, Statistical, and Arithmetic Space}
\label{sec:interpretation}

This section is conceptual, not a theorem statement, and should be read as orientation for a reader extending this work rather than as a claim.

It is useful to separate three levels at which a parity word can be described:

\begin{description}[leftmargin=2em]
\item[Symbolic space.] The actual word: its literal order, block structure, and factor-by-factor arrangement.
\item[Statistical projection.] Aggregated summaries of the symbolic object --- density $\gamma=k/\ell$, discrepancy $D(W)$, block frequencies, factor complexity. Many distinct symbolic objects can share the same statistical projection.
\item[Arithmetic admissibility.] Which symbolic objects, specifically, correspond to a rational $\Phi$-value. This is a pointwise, not a statistical, property.
\end{description}

The relationship is a \emph{fiber}: symbolic object $\to$ statistical state is generally many-to-one. A crucial consequence, easy to state and easy to get wrong in either direction: \textbf{statistical sparsity of a class (low complexity, small discrepancy, controlled density) does not by itself imply arithmetic impossibility of rational $\Phi$-values within it}, and conversely, statistical richness does not by itself force irrationality. The bridge proved in this note works precisely because it does \emph{not} stop at a statistical argument: it lifts symbolic and statistical information (repetition length, discrepancy) back into a pointwise arithmetic contradiction (Theorem \ref{thm:abstract}'s exact integer divisibility argument) for the \emph{specific} word $s$ under study, not for a statistical ensemble containing it. This is why the method scales from Sturmian words to CS Rote sequences without a change in kind --- and also why it stalls exactly where the statistical control (discrepancy, excess) it depends on stalls (\S\ref{sec:generalintercept}).

\section{The Drift/Resonance Coordinate}
\label{sec:drift}

For a finite word $W$ ($\ell=|W|$, $k=|W|_1$), define
\[
R(W) := \ell - k\log_2 3.
\]
Then $2^\ell - 3^k = 2^\ell\bigl(1-2^{-R(W)}\bigr)$, and near $R=0$ (i.e.\ near the ``resonance'' $\ell\approx k\log_2 3$), $|2^\ell-3^k| \approx 2^\ell \ln2\,|R(W)|$.

This coordinate is explanatory, not structural: it was investigated directly in Phase 7 as a candidate new invariant, and the record is explicit that it is not one. \emph{``No new information is encoded by $R$ that $\gamma$ did not already carry''} \cite{Phase7Drift}; \emph{``Resonance did not materially help or hurt in any tested case''} \cite{Phase7DriftDigits}. $R(W)$ is a reparameterization of the density margin $\gamma\log_2 3$ already used throughout \S\ref{sec:repheight}--\S\ref{sec:oddweight}, useful for visualizing how close a given witness sits to the arithmetic boundary, but supplying no additional proof-theoretic leverage found anywhere in this program.

\section{The Symmetry-Tradeoff Probe (Exploratory, Inconclusive)}
\label{sec:symmetryprobe}

\textbf{This section reports an exploratory investigation. Nothing in it is a theorem, and nothing in it should be read as evidence for or against any conjecture stated elsewhere in this note.} It is included so a future researcher does not repeat the same probe without knowing what was already tried and why it did not resolve anything.

The idea investigated: does \emph{approach to symmetry} in the Ostrowski digit sequence (roughly, how often the maximal admissible continuation is available) trade off against arithmetic \emph{resonance} (how close $R(W)$, \S\ref{sec:drift}, sits to $0$)? Two experiments were run.

\textbf{Experiment 1} correlated a normalized arithmetic-slack quantity $A(W):=|1-2^{-R(W)}|$ against the exact surplus $S(W)$ and against the achieved repetition ratio $L/\ell$, using $74$ genuine witnesses from the general-intercept search of \S\ref{sec:generalintercept}. Result: $A(W)$ vs.\ $S(W)$, $r\approx0.208$; $A(W)$ vs.\ $L/\ell$, $r\approx-0.125$. Both fall below the probe's own $0.3$ threshold for ``worth a theorem-seeking follow-up,'' and $n=74$ is too small for $r=0.208$ to clear the standard significance threshold for that sample size. Separately, the first correlation is only weak evidence in any case: the surplus machinery already leans on the same $\ell-k\log_2 3$ resonance quantity that defines $A(W)$, so a correlation there may be partly definitional rather than an independent discovery.

\textbf{Experiment 2}, testing a genuinely independent quantity (a windowed frequency of admissible-digit freedom, $N_{\text{sym}}$, against $A(W)$), did not run to completion in its original form: it materialized an explicit substitution-image word whose length grows exponentially in the search depth $k$ (for slope $A=6$, a projected $\sim9.9\times10^{14}$ symbols by $k=19$), exhausting memory. This was diagnosed and repaired: the word's length and weight can be obtained exactly, without ever materializing it, via a $2\times2$ integer transfer-matrix representation of the same substitutions ($\tau_0,\tau_1$; see \S\ref{sec:transfermatrix}), cross-checked against brute-force construction in $150$ of $150$ small cases with exact agreement. With this fix, Experiment 2 completed on $724$ points: $A(W)$ vs.\ $N_{\text{sym}}$, $r\approx-0.0058$; on a log--log scale, $r\approx-0.0319$. \textbf{These numbers must not be read at face value}: $596$ of the $724$ points ($82\%$) had $A(W)$ pinned to exactly $1.0$ by the probe's own floating-point underflow clamp ($|R(W)|>60$), an artifact of the diagnostic's float arithmetic, not a genuine measurement in the near-resonance regime where any real effect would have to live.

\begin{quote}
\textbf{Verdict: INCONCLUSIVE.} Neither experiment produces a correlation free of confounding structure (definitional entanglement in Experiment 1; a numerical clamping artifact dominating Experiment 2). This is consistent with, and independently corroborates, the Phase 7 drift-coordinate finding of \S\ref{sec:drift} that resonance carries no separately exploitable signal in every case tested so far --- but it establishes nothing new, and \emph{the symbolic-symmetry/arithmetic-freedom tradeoff hypothesis is neither supported nor contradicted} by this probe.
\end{quote}

A methodologically sound re-run would need $A(W)$ computed at sufficient precision to avoid the clamp, restricted to (or stratified by) genuine near-resonance witnesses, before any correlation from it could be informative.

\section{A Transfer-Matrix Computational Tool}
\label{sec:transfermatrix}

\textbf{This section records a computational/structural tool, not a new irrationality theorem.} It is preserved here because it is broadly useful and easy to lose track of.

The Sturmian substitutions used throughout this program (e.g., in constructing explicit BHZ witness words, \S\ref{sec:generalintercept}) are
\[
\tau_0: 0\mapsto 0,\ 1\mapsto 01; \qquad \tau_1: 0\mapsto 10,\ 1\mapsto 1.
\]
Applying either substitution repeatedly can make the resulting word's length grow exponentially in the number of applications --- and that length is all that is needed for most downstream computations (arithmetic slack $A(W)$, weight $\kappa$, density). It need never be materialized. Tracking only $(n_0,n_1)$, the counts of $0$s and $1$s, the two substitutions act linearly:
\[
\tau_0: (n_0,n_1)\mapsto(n_0+n_1,\ n_1), \qquad
\tau_1: (n_0,n_1)\mapsto(n_0,\ n_0+n_1),
\]
i.e.\ by the integer matrices $M_0=\begin{psmallmatrix}1&1\\0&1\end{psmallmatrix}$, $M_1=\begin{psmallmatrix}1&0\\1&1\end{psmallmatrix}$ (both determinant $1$; their products are exactly the continuant recurrence generating the convergent denominators $q_k$). Applying $\tau_i$ exactly $e$ times in a row is $M_i^e$; composing across an entire substitution schedule of length $k$ is $k$ integer matrix--vector multiplications, giving the word's exact length $\ell=n_0+n_1$ and weight $\kappa=n_1$ in $O(k)$ exact-integer steps regardless of how astronomically large $\ell$ becomes --- verified against brute-force construction with $150/150$ exact agreement wherever brute force was still feasible to run.

A related, separate observation: counting the number of admissible Ostrowski digit sequences $(c_k)$ (rather than computing one fixed sequence's word) is itself governed by a minimal two-state automaton, not by any larger state. The admissibility rule (\S\ref{sec:generalintercept}) depends only on whether the immediately preceding digit was exactly $0$: from that state, the next digit has $a_k+1$ admissible choices; from any nonzero-previous state, only $a_k$. No memory of the actual previous digit value, nor of anything beyond the current index into $(a_k)$, is needed. This directly mirrors the admissible-sequence generator already used throughout the repository's scripts, and is recorded here as a structural fact about that generator, not as a separately theorem-proving result.

\section{What Is Still Open}
\label{sec:open}

\begin{description}[leftmargin=3em,style=nextline]
\item[Open 1 --- General intercept margin. \textnormal{\textbf{CLOSED} (Revision 2, C1).}] Revision
1's Conjecture 12.1 is BHZ Propositions 5.1 and 5.2 (Theorem \ref{thm:bhzmargin}), with an
attained-witness refinement in Theorem \ref{thm:generalmargin}. Theorem \ref{thm:boundedtype} now
holds at every intercept. What remains here is only a sharpness question, not an obstruction: is
$\frac1{A\beta_A}$ the infimum of $\text{ice}-2$ on slopes whose digit tail is \emph{not} eventually
constant? Proved for eventually-constant tails with $A\ge2$; numerical elsewhere
\cite{Phase8Margin}. Nothing depends on the answer.
\item[Open 2 --- Return/exceptional intercept set. \textnormal{\textbf{WITHDRAWN} (Revision 2, C2).}]
The premise was false: the set of intercepts with $\text{ice}(\omega)<\text{ind}^*(\gamma)$ is
uncountable, not a countable return set, and intercept $0$ sits among the \emph{easiest} intercepts
rather than the hardest. There is no distinguished countable family left to classify.
\item[Open 3 --- Intrinsic-square bypass. \textnormal{\textbf{SUPERSEDED} (reading confirmed).}] The
Revision 1 analysis was correct --- bare squares ($r=2$) give a fixed margin $2-\log_23$ on the
even-weight branch and \emph{exactly zero} on the odd-weight branch \cite{Phase7Intrinsic} --- and
Theorem \ref{thm:generalmargin} makes the bypass unnecessary by supplying $r>2$ strictly at every
intercept. The route retains one distinct virtue, recorded in Open 4: Allouche--Davison--Queff\'elec--Zamboni's
square theorem has \emph{no} hypothesis on $\gamma$, so it is the only part of the chain that
survives unbounded partial quotients.
\item[Open 4 --- Unbounded partial quotients. \textnormal{\textbf{THE LIVE FRONTIER}.}] Remove the
bounded-partial-quotient hypothesis from Theorem \ref{thm:boundedtype}. This is no longer merely
``flagged as hard'': BHZ Theorem 1.1 \emph{characterises} the slopes carrying an intercept with
$\text{ice}(\omega)=2$ exactly --- for each pair $(s,t)$, $s>1$, only finitely many $k$ with
$(a_k,a_{k+1})=(s,t)$ or $(a_k,a_{k+1},a_{k+2})=(1,1,t)$ --- and every such slope has unbounded
partial quotients. At such an intercept the odd-weight branch has margin exactly $0$, so the present
sufficient criterion cannot apply, independently of any discrepancy estimate. Three distinct
questions, in increasing order of ambition: (i) can the \emph{even} branch be forced infinitely
often, where bare squares already give $2-\log_23$ and the discrepancy hypothesis can be dropped
entirely? (ii) is the $O(\log\ell)$ discrepancy bound replaceable by an exact height estimate that
tolerates $E_k=0$? (iii) does a CS Rote sequence carry intrinsic initial repetition of its own,
beyond what transfers from $u$? See \repo{theorems/phase8/PHASE8\_NEXT\_TARGET.md}.
\item[Open 5 --- Higher complexity.] Only once the $p(n)=2n$ rung is fully closed (Open 1--4): which symbolic classes beyond Rote sequences --- higher, structured complexity classes; general two-interval rotation codings, for which the repetition-existence half of the argument is itself still open \cite{Phase2Rotation}; general morphic words --- can this bridge reach next?
\item[Open 6 --- The full Periodicity Conjecture.] What mechanism, if any, could cover aperiodic words that admit no useful initial repetition structure at all? The Thue--Morse word has no square whatsoever below half-length $2048$ \cite{Phase1Gap}, making every method in this note --- indeed, every purely repetition-based method --- structurally inapplicable to it by construction. This is not a gap to be closed by extending the present technique further; it is evidence that a fundamentally different mechanism would be needed for the conjecture in full generality.
\end{description}

\section{Future-Work Roadmap}
\label{sec:roadmap}

The roadmap below is organized directly around the six open problems of \S\ref{sec:open} and the repository's own stated priorities.

\begin{description}[leftmargin=3em,style=nextline]
\item[Near term.] Open 1--3 are settled or withdrawn, so the near-term work is not mathematical but
editorial and hygienic: propagate the $q_1=a_1+1$ convention through the underlying record's code
and re-derive every root length reported there (no $\text{ice}$ value changes), and fold Lemma
\ref{lem:halfslope} into the discrepancy audit where it was previously implicit.
\item[Medium term.] Open 4(i): a general-intercept weight-parity criterion. The even-weight branch
survives $\text{ice}(\omega)=2$ with margin $2-\log_23\approx0.415$ and needs no discrepancy
hypothesis, so it suffices to force even weight infinitely often. At intercept $0$ this was the
Phase 3 automaton; at general intercepts the prefix-count parity is the classical Ostrowski-digit
sum $k_i(u)=\lfloor i\gamma+\rho\rfloor-\lfloor\rho\rfloor$ read mod $2$, so a literature search on
Ostrowski-digit parity should precede any new derivation. This is the single highest-value target
identified by Revision 2.
\item[Longer term.] Open 4(ii)--(iii): an exact height estimate tolerating zero repetition excess,
or intrinsic Rote repetition not inherited from $u$. Both would be genuinely new mechanisms rather
than extensions, and both should be attempted only after 4(i) is decided.
\item[Structural.] Any attempt at Open 6 (the full conjecture) needs a mechanism that does not depend on initial repetition existing at all, given the Thue--Morse-type structural ceiling recorded in \S\ref{sec:open}. This program takes no position on what such a mechanism could look like; it records only that repetition-based methods, however far extended, cannot be it alone.
\item[Process.] Continue the falsification-first discipline of \S\ref{sec:selfcorrection}: every extension of this program should include an adversarial audit pass against genuinely non-periodic, non-quadratic controls before any theorem upgrade is accepted, exactly as Phases 5 and 6 did for the bounded-type theorem.
\end{description}

\section*{Provenance and Verification}
\addcontentsline{toc}{section}{Provenance and Verification}

\textbf{Revision 2} was prepared from the \repo{periodicity-conjecture-bridge} repository on the
isolated branch \repo{phase8-general-intercept}, whose Phase 8 deliverables
(\repo{theorems/phase8/}, \repo{scripts/phase8/}, \repo{data/phase8/}) carry every corrected
statement, and which was not pushed during this preparation. No Phase 1--7 research file was
rewritten: the Phase 7 record is preserved verbatim and carries conspicuous correction banners
pointing to \repo{theorems/phase7/PHASE7\_CORRECTIONS.md} \cite{Phase7Corrections}. Revision 1 was consolidated from the same
repository at commit \texttt{a4ac46aebf17fb2ce14cea50f5ce60a58704da15} (``Phase 7: map
general-intercept CS Rote obstruction''). Every theorem statement, definition, and quoted result above was extracted directly from the corresponding file under \repo{theorems/phase\{1..7\}/} and cross-checked against \repo{THEOREM\_STATUS.md} and \repo{LADDER.md}; none was reconstructed from memory. SHA-256 checksums of every file under \repo{theorems/}, \repo{scripts/}, \repo{data/}, and \repo{reference/} at the time of writing are recorded in \repo{technical-note/PRESERVATION\allowbreak\_RECORD.md}. A full statement-by-statement classification (known/cited, proved in foundation paper, proved in this program, computationally verified, counterexample, rejected hypothesis, open, or exploratory only) is recorded in \repo{technical-note/SOURCE\_OF\_TRUTH.md}. No Phase 1--7 research file was modified in the preparation of this note.

Revision 1 is deposited as Zenodo version DOI \texttt{10.5281/zenodo.23045141}; Revision 2 is
intended as a \emph{new version} of that record, leaving Revision 1 retrievable, not as an edit of
it. A concise standalone erratum accompanies it (\repo{technical-note/correction\_notice.tex}).

\section*{Acknowledgments}
\addcontentsline{toc}{section}{Acknowledgments}
AI assistance (Anthropic's Claude) was used for repository archaeology, adversarial diagnostic work (including diagnosing and repairing the Experiment 2 performance bug in \S\ref{sec:symmetryprobe}), literature cross-checking against the foundation paper's own bibliography, and drafting of this note, consistent with the acknowledgment practice of the foundation paper \cite{DeJesus2026Sturmian}. All mathematical claims, their framing, and any errors are the author's.

\begin{thebibliography}{99}

\bibitem{DeJesus2026Sturmian} Elias De Jesús, \emph{The 3x+1 Conjugacy Map Sends Every Sturmian Word to an Irrational 2-Adic Integer}, unpublished manuscript, September 2026. (Foundation paper for this research program; companion work by the present author.)

\bibitem{ADQZ2001} Jean-Paul Allouche, J.\ L.\ Davison, Michel Queffélec, and Luca Q.\ Zamboni, \emph{Transcendence of Sturmian or morphic continued fractions}, Journal of Number Theory \textbf{91} (2001), no.\ 1, 39--66.

\bibitem{Bernstein1994} Daniel J.\ Bernstein, \emph{A noniterative 2-adic statement of the $3N+1$ conjecture}, Proceedings of the American Mathematical Society \textbf{121} (1994), no.\ 2, 405--408.

\bibitem{BernsteinLagarias1996} Daniel J.\ Bernstein and Jeffrey C.\ Lagarias, \emph{The $3x+1$ conjugacy map}, Canadian Journal of Mathematics \textbf{48} (1996), no.\ 6, 1154--1169.

\bibitem{BHZ2006} Valérie Berthé, Charles Holton, and Luca Q.\ Zamboni, \emph{Initial powers of Sturmian sequences}, Acta Arithmetica \textbf{122} (2006), no.\ 4, 315--347.

\bibitem{Dubickas2009} Artūras Dubickas, \emph{On integer sequences generated by linear maps}, Glasgow Mathematical Journal \textbf{51} (2009), no.\ 2, 243--252.

\bibitem{Lagarias1985} Jeffrey C.\ Lagarias, \emph{The $3x+1$ problem and its generalizations}, The American Mathematical Monthly \textbf{92} (1985), no.\ 1, 3--23.

\bibitem{LopezStoll2009} Josefina López and Peter Stoll, \emph{The $3x+1$ conjugacy map over a Sturmian word}, Integers \textbf{9} (2009), no.\ A13, 141--162.

\bibitem{LopezStoll2021} Josefina López and Peter Stoll, \emph{The $3x+1$ periodicity conjecture in $\mathbb{R}$}, 2021.

\bibitem{Lothaire2002} M.\ Lothaire, \emph{Algebraic Combinatorics on Words}, Encyclopedia of Mathematics and its Applications, vol.\ 90, Cambridge University Press, 2002, Chapter 2 (Sturmian words, by J.\ Berstel and P.\ Séébold). Cited here as the standard reference for the combinatorics-on-words background of complementary symmetric Rote sequences; the primary source defining this class (G.\ Rote, 1994) should be verified directly against primary literature before formal submission, as this repository's record does not independently confirm its exact publication details.

\bibitem{Phase2Abstract} \repo{theorems/phase2/ABSTRACT\_IRRATIONALITY\_CRITERION.md}, periodicity-conjecture-bridge repository, commit \texttt{a4ac46a}.
\bibitem{Phase2Surplus} \repo{theorems/phase2/SURPLUS\_THEOREM.md} and \repo{theorems/phase1/SURPLUS\_INVARIANT.md}, ibid.
\bibitem{Phase1Discrepancy} \repo{theorems/phase1/DISCREPANCY\_HEIGHT\_THEOREM.md}, ibid.
\bibitem{Phase1Repetition} \repo{theorems/phase1/REPETITION\_HEIGHT\_CRITERION.md}, ibid.
\bibitem{Phase2Discrepancy} \repo{theorems/phase2/ROTE\_ROOT\_DISCREPANCY.md}, ibid.
\bibitem{Phase2Transfer} \repo{theorems/phase2/ROTE\_TRANSFER\_THEOREM.md}, ibid.
\bibitem{Phase5Transfer} \repo{theorems/phase5/TRANSFER\_REDERIVATION.md}, ibid.
\bibitem{Phase3Automaton} \repo{theorems/phase3/WEIGHT\_PARITY\_CLASSIFICATION.md}, ibid.
\bibitem{Phase4Silver} \repo{theorems/phase4/SILVER\_EXTREMALITY.md}, ibid.
\bibitem{Phase4SilverAudit} \repo{theorems/phase4/SILVER\_CONSTANT\_AUDIT.md}, ibid.
\bibitem{Phase4Required} \repo{theorems/phase4/REQUIRED\_EXPONENT.md}, ibid.
\bibitem{Phase5Hostile} \repo{theorems/phase5/HOSTILE\_REFEREE\_REPORT.md}, ibid.
\bibitem{Phase6Arrow} \repo{theorems/phase6/ARROW\_BY\_ARROW\_RECONSTRUCTION.md}, ibid.
\bibitem{Phase4Quadratic} \repo{theorems/phase4/QUADRATIC\_CS\_ROTE\_THEOREM.md}, ibid.
\bibitem{Phase6Bounded} \repo{theorems/phase6/BOUNDED\_TYPE\_THEOREM.md}, ibid.
\bibitem{Phase4Chain} \repo{theorems/phase4/FINAL\_PROOF\_CHAIN.md}, ibid.
\bibitem{BHZGeneralMap} \repo{theorems/phase7/BHZ\_GENERAL\_INTERCEPT\_MAP.md}, ibid.
\bibitem{Phase7Verdict} \repo{theorems/phase7/PHASE7\_VERDICT.md}, ibid.
\bibitem{Phase7Sufficient} \repo{theorems/phase7/INTERCEPT\_SUFFICIENT\_CONDITION.md}, ibid.
\bibitem{Phase7Surplus} \repo{theorems/phase7/SURPLUS\_DECOMPOSITION.md} and \repo{theorems/phase7/INTERCEPT\allowbreak\_EXCESS.md}, ibid.
\bibitem{Phase7Flexible} \repo{theorems/phase7/FLEXIBLE\_MARGIN\_BOUNDARY.md}, ibid.
\bibitem{Phase7Structured} \repo{theorems/phase7/STRUCTURED\_INTERCEPTS.md}, ibid.
\bibitem{Phase7Drift} \repo{theorems/phase7/DRIFT\_COORDINATES.md}, ibid.
\bibitem{Phase7DriftDigits} \repo{theorems/phase7/DRIFT\_VS\_DIGITS.md}, ibid.
\bibitem{Phase7Intrinsic} \repo{theorems/phase7/INTRINSIC\_SQUARE\_ROUTE.md}, ibid.
\bibitem{Phase2Rotation} \repo{theorems/phase2/ROTATION\_EXAMPLE\_THEOREM.md}, ibid.
\bibitem{Phase1Gap} \repo{theorems/phase1/FULL\_PERIODICITY\_GAP.md}, ibid.

\bibitem{Phase8Proof} \repo{theorems/phase8/PHASE8\_PROOF\_OR\_OBSTRUCTION.md}, periodicity-conjecture-bridge repository, branch \repo{phase8-general-intercept}.
\bibitem{Phase8Target} \repo{theorems/phase8/PHASE8\_TARGET\_AUDIT.md}, ibid.
\bibitem{Phase8Chain} \repo{theorems/phase8/PHASE8\_CHAIN\_AUDIT.md}, ibid.
\bibitem{Phase8Margin} \repo{theorems/phase8/PHASE8\_MARGIN\_STATUS.md}, ibid.
\bibitem{Phase8Verdict} \repo{theorems/phase8/PHASE8\_VERDICT.md}, ibid.
\bibitem{Phase7Corrections} \repo{theorems/phase7/PHASE7\_CORRECTIONS.md}, ibid.

\end{thebibliography}

\end{document}
